\chapter{Conclusiones y Trabajo Futuro}
\label{cap:conclusiones}

\section{Conclusiones}

\subsection{Sobre las técnicas de adaptación}

De las tres técnicas evaluadas, solo mejora la que modifica los pesos del modelo. El
ajuste fino con LoRA reduce el error 1,23 puntos con ambos criterios estadísticos de
acuerdo; el prompting contextual no produce efecto detectable, con signo distinto en cada
corpus y pruebas que discrepan entre sí; y la post-corrección con un modelo de lenguaje
local lo empeora de forma significativa y replicada.

La conclusión sobre la post-corrección merece detenerse, porque contradice la intuición
del campo. Un modelo de lenguaje debería poder reparar errores de reconocimiento usando el
contexto, y sin embargo degrada. El mecanismo observado explica por qué: el corrector
optimiza la calidad del texto como pieza escrita, sustituyendo palabras por sinónimos más
precisos y corrigiendo la gramática del hablante, mientras que una transcripción debe
reflejar lo dicho. Son objetivos en conflicto, y el modelo optimiza el que le corresponde
por entrenamiento.

Sobre LoRA cabe un matiz que conviene no omitir: buena parte de su mejora consiste en
aprender la convención de escritura del corpus de ajuste, no en reconocer mejor el audio.
Descontado ese efecto, la ganancia sigue existiendo pero es modesta.

\subsection{Sobre dónde está realmente el margen de mejora}

El resultado más relevante del trabajo no procede de la comparativa de técnicas sino de su
comparación con las decisiones de ingeniería del sistema. Segmentar el audio por silencios
en lugar de por reloj reduce el error 7,4 puntos a igual latencia media; seis veces más que
la mejor de las técnicas de adaptación, y a coste de cómputo prácticamente nulo.

Esto invierte la premisa de la que partía el trabajo. Se asumía que el margen estaba en
adaptar el modelo al dominio educativo; los datos indican que, con los modelos actuales, el
margen está en cómo se le entrega el audio. Un sistema de subtitulado en vivo que trocee
mal desperdicia más calidad de la que cualquier técnica de adaptación puede recuperar.

\subsection{Sobre la evaluación}

La tasa de error de palabra es insuficiente para valorar un sistema de accesibilidad, y el
trabajo aporta cuatro evidencias independientes de ello. La más contundente: con un 18,61\%
de error, cifra que suena aceptable, el sistema pierde o altera el 28\% de las negaciones e
introduce trece que nadie pronunció. Un alumno con discapacidad auditiva lee entonces lo
contrario de lo explicado, sin ninguna señal de que algo falla.

De ahí se derivan dos instrumentos que se proponen como complemento: una métrica de
cobertura y precisión sobre la terminología del dominio, y otra que contabiliza los errores
sobre negaciones, numerales y cuantificadores. Ambas detectan efectos que el WER agregado
no muestra, y en un caso concreto invierten la valoración de una técnica.

\subsection{Sobre la validez de las mediciones}

Cinco incidentes documentados en la sección~\ref{sec:incidentes} comparten una
característica: en todos ellos el sistema seguía funcionando y produciendo cifras
verosímiles. Ninguno se manifestó como un error.

La conclusión práctica es que un procedimiento experimental en este campo necesita
comprobaciones que no forman parte del procedimiento habitual: verificar que la función de
pérdida esté en el rango esperado antes de evaluar un modelo ajustado; auditar a mano una
muestra de las referencias antes de adoptar un corpus; incorporar la configuración de
decodificación a la identidad del resultado; y registrar el contenido exacto del corpus
empleado, no solo su nombre.

\section{Limitaciones}
\label{sec:limitaciones}

\textbf{Ninguno de los corpus es una clase real.} Es la limitación de fondo. Los materiales
empleados se aproximan al registro de una clase magistral, pero ninguno lo es. Las
conclusiones sobre adaptación al dominio educativo son, en rigor, conclusiones sobre
adaptación a los dominios sustitutos disponibles. El corpus objetivo, un repositorio
universitario con transcripciones manuales, requiere una solicitud institucional que quedó
cursada sin resolución dentro del plazo de este trabajo.

\textbf{Los resultados miden implementaciones, no técnicas.} Se hizo evidente con la
post-corrección: una primera versión de las instrucciones dejaba la técnica inerte, y otra
la activaba. Un prompt distinto, un modelo mayor o un ajuste de hiperparámetros podrían
alterar las conclusiones.

\textbf{El detector de actividad de voz es de energía.} Un detector entrenado sería más
robusto frente al ruido de aula, y la práctica del sector indica que la combinación de
ambos reduce sustancialmente las falsas activaciones. La medición realizada justifica
introducirlo, pero no se ha implementado.

\textbf{El ruido se añadió de forma sintética.} Falta la reverberación, que es un factor
independiente y posiblemente más dañino que el ruido aditivo en un aula real.

\textbf{La aplicación no se ha evaluado con usuarios.} Se ha verificado su funcionamiento
técnico y se ha medido su latencia, pero no se ha comprobado si el resultado es útil para
quien debe leerlo.

\textbf{El sistema depende de un único servidor.} La difusión se resuelve en memoria, lo
que ata el despliegue a un proceso. Escalar a varios exigiría una infraestructura
adicional; para un aula no aplica, pero condiciona cualquier extensión a nivel de centro.

\section{Trabajo futuro}

\textbf{Corpus educativo real.} Es la continuación prioritaria: repetir la comparativa
sobre clases grabadas convertiría las conclusiones actuales, obtenidas sobre dominios
sustitutos, en conclusiones sobre el dominio objetivo.

\textbf{Detección de actividad de voz entrenada.} La segmentación resultó ser la palanca
más eficaz del sistema, y actualmente se apoya en el método más simple posible. Mejorarla
es la vía con mayor rendimiento esperable por unidad de esfuerzo.

\textbf{Glosario procedente del material docente.} La extracción automática por frecuencia
produce vocabulario corriente que el modelo ya reconoce. Un glosario elaborado a partir del
temario o las transparencias tendría el contenido que la técnica necesita, y permitiría
reevaluar el prompting en condiciones justas.

\textbf{Métricas de comprensión.} Las métricas propuestas miden la supervivencia de
unidades léxicas concretas. El paso siguiente sería medir si el alumno entiende lo mismo
leyendo la transcripción que escuchando la clase, lo que exige evaluación con personas.

\textbf{Implementaciones optimizadas para inferencia.} Existen variantes del modelo
significativamente más rápidas que la empleada, cuyo soporte para la plataforma de cálculo
utilizada en este trabajo era limitado. Su viabilidad debería reevaluarse, porque el margen
de latencia condiciona directamente la usabilidad en aula.
